\section{Edge vs. Ground Processing}

\subsection{Experimental Setup}

This section evaluates the trade off between running human detection on-board each UAV (edge) and running it at the ground station (ground). The two configurations use an identical detector YOLOv8n restricted to the COCO person class at a confidence threshold of 0.25 and differ only in where inference is placed and, consequently what crosses the the simulated link in edge case is compact detection messages and in ground case it's JPEG compressed imagery (quality 5). Detection accuracy was assessed on $60$ manually labelled frames, evaluated in both representations (raw for edge, JPEG quality 5 for ground) so that the comparison is done on identical conditions. Experiments ran on the simulation host (Intel Core i7-13700, 30GiB RAM, Ubuntu 22.04) with ROS2 Humble, Gazebo and ns-3.38. The UAV flew at 25m relative altitude with the cameral relayed at 1Hz over a 60s measurement window.

Each configuration was evaluated using three independent ns3 RNG runs (seeds 1-3). One complete run was treated as a single expremental repetition where individual frames were not treated as independent samples, since frames within a run are correlated. All collected results are therefore reported as mean $\pm$ sample standard deviation across the three runs. To exclude transient start-up behaviour, the official window of each run began only after a settling period and a stability criterion of several consecutive stable frame arrivals were met. So that the measurement covered steady state operation rather than pipeline warm up. 

Four categories of metrics were recorded. Detection accuracy is reported as precision, recall, F1-score and exact count rate (the fraction of frames on which the detected person count exactly matched the label). Processed frame ratio measures pipeline reliability. The proportion of frames published by the camera relay during the official that produced a corresponding detection result within that window. Pipeline latency is roth configurations were executed and monitored on the simulation host. The edge configuration is additionally realised on physical Raspberry Pi hardware-in-the-loop discussed in Section VI and the host based resource figures here are not a substitute for on device power measurement. 


\subsection{Detection Accuracy}



\subsection{Reliability and Latency}



\subsection{Resource Utilisation}


